\section{محاسبۀ انتگرال‌های معین به روش مانده‌ها}
\def\EMfootchapter{فصل ۱۴ \ \lr{\textbullet}\ محاسبۀ انتگرال‌های معین به روش مانده‌ها}
\def\EMtitlesize{26}\def\EMtitleskip{31}
\begin{frame}{فصل ۱۴: محاسبۀ انتگرال‌های معین به روش مانده‌ها}
{\fontsize{20}{26}\selectfont\bfseries\color{SLdef}\EMalign مطالب این فصل\par}
\vspace{10bp}
\EMitem{\lr{\textbullet}}{\dimexpr\linewidth-26bp\relax}{انتگرال‌های دستۀ اول}
\EMitem{\lr{\textbullet}}{\dimexpr\linewidth-26bp\relax}{انتگرال‌های دستۀ دوم}
\end{frame}

\def\EMtitlesize{26}\def\EMtitleskip{31}
\begin{frame}{محاسبۀ انتگرال‌های معین به روش مانده‌ها}
\begin{EMbox}{thm}{1}{}
\begin{uncoverenv}<2->
\EMp{اگر تابع \EMim{f(z)} در همه‌جای داخل مسیر \EMim{C} به‌جز نقاط \EMim{z_1,z_2,\dots,z_n} تحلیلی باشد و مانده تابع در این نقاط به ترتیب \EMim{A_1^1,A_1^2,\dots,A_1^n} باشد، آنگاه:}
\vspace{3bp}
\EMeq{\oint_Cf(z)\,dz =2\pi i\sum_{k=1}^{n}A_1^k}
\end{uncoverenv}
\end{EMbox}
\begin{uncoverenv}<3->
\EMp{دسته انتگرال‌های مختلفی توسط روش مانده‌ها قابل حل هستند. در اینجا به دو دستۀ آنها اشاره می‌شود:}
\end{uncoverenv}
\EMgap
\begin{EMbox}{def}{3}{}
\begin{uncoverenv}<4->
\EMp{در این نوع انتگرال‌ها، مخرج انتگرند به صورت چندجمله‌ای بوده، صورت به فرم تابعی از نوع چندجمله‌ای یا نمائی و یا توابع سینوس یا کسینوس (که آنها هم به صورت نمائی قابل بیان هستند) می‌باشد.}
\end{uncoverenv}
\end{EMbox}
\begin{EMbox}{def}{4}{}
\begin{uncoverenv}<5->
\EMp{در این نوع انتگرال‌ها، انتگرند تابعی از سینوس یا کسینوس می‌باشد.}
\end{uncoverenv}
\end{EMbox}
\end{frame}

\def\EMtitlesize{26}\def\EMtitleskip{31}
\begin{frame}{انتگرال‌های دستۀ اول}
\begin{EMbox}{ex}{1}{}
\EMp{مطلوب است محاسبۀ انتگرال \EMim{\displaystyle\int_0^\infty\frac{\cos\omega x}{(1+x^2)^2}\,dx} (\EMim{\omega>0}).}
\EMgap
\begin{uncoverenv}<2->
\EMp{تابع \EMim{f(z)=\dfrac{e^{i\omega z}}{(1+z^2)^2}} را در نظر می‌گیریم. این تابع در \EMim{z=i} و \EMim{z=-i} تحلیلی نیست. مسیر \EMim{C} نیم‌دایرۀ بالایی به شعاع \EMim{R} است.}
\vspace{3bp}
\EMeq{\begin{aligned}
\oint_C\frac{e^{i\omega z}}{(1+z^2)^2}\,dz&=2\pi i\times\operatorname{Res}\left\{\frac{e^{i\omega z}}{(1+z^2)^2}\right\}\Bigg|_{z=i}\\
&=2\pi i\times\left[\frac d{dz}\left\{(z-i)^2\,\frac{e^{i\omega z}}{(z-i)^2(z+i)^2}\right\}\right]_{z=i}\\
&=2\pi i\times\left[\frac{i\omega e^{i\omega z}(z+i)-2e^{i\omega z}}{(z+i)^3}\right]_{z=i}\\
&=2\pi i\times e^{-\omega}\left(\frac{i\omega(2i)-2}{(2i)^3}\right)\\
&=2\pi i\times e^{-\omega}\,\frac{\omega+1}{4i}=\frac\pi2(\omega+1)e^{-\omega}
\end{aligned}}
\end{uncoverenv}
\end{EMbox}
\end{frame}

\def\EMtitlesize{26}\def\EMtitleskip{31}
\begin{frame}{انتگرال‌های دستۀ اول (ادامه)}
\begin{minipage}[t]{500bp}\vspace{0pt}
\begin{EMbox}{ex}{1}{}
\begin{uncoverenv}<2->
\EMp{از طرف دیگر}
\vspace{3bp}
\EMeq{\oint_C\frac{e^{i\omega z}}{(1+z^2)^2}\,dz \\[4pt] =\int_{-R}^{R}\frac{e^{i\omega x}}{(1+x^2)^2}\,dx+\int_0^\pi\frac{e^{i\omega Re^{it}}\,Rie^{it}}{(1+R^2e^{i2t})^2}\,dt}
\end{uncoverenv}
\EMgap
\begin{uncoverenv}<3->
\EMp{و وقتی \EMim{R\to\infty} انتگرال روی نیم‌دایره صفر می‌شود و بخش موهومی انتگرال روی محور حقیقی (تابع فرد \EMim{\dfrac{\sin\omega x}{(1+x^2)^2}}) صفر است:}
\vspace{3bp}
\EMeq{\oint_C\frac{e^{i\omega z}}{(1+z^2)^2}\,dz\ \xrightarrow{\ R\to\infty\ }\ 2\int_0^\infty\frac{\cos\omega x}{(1+x^2)^2}\,dx}
\end{uncoverenv}
\EMgap
\begin{uncoverenv}<4->
\EMeq{\int_0^\infty\frac{\cos\omega x}{(1+x^2)^2}\,dx =\frac\pi4(\omega+1)e^{-\omega}}
\end{uncoverenv}
\end{EMbox}
\end{minipage}
\EMside{136.5bp}{\uncover<5->{\EMfigure{1.500}{m14-semicircle}{مسیر نیم‌دایرۀ بالایی؛ قطب مرتبۀ دوم \EMim{z=i} داخل مسیر و \EMim{z=-i} بیرون آن است}}}
\end{frame}

\def\EMtitlesize{26}\def\EMtitleskip{31}
\begin{frame}{انتگرال‌های دستۀ دوم}
\begin{uncoverenv}<2->
\EMeq{\int_0^{2\pi}f(a\cos t,\,b\sin t)\,dt}
\end{uncoverenv}
\EMgap
\begin{uncoverenv}<3->
\EMp{برای حل این‌گونه انتگرال‌ها فرض می‌کنیم \EMim{z=e^{it}}، در این صورت داریم:}
\vspace{3bp}
\EMeq{dz =ie^{it}\,dt\; \Longrightarrow\;dt =\frac{dz}{ie^{it}} =\frac{dz}{iz}}
\end{uncoverenv}
\EMgap
\begin{uncoverenv}<4->
\EMeq{\sin t =\frac{e^{it}-e^{-it}}{2i} =\frac{z-z^{-1}}{2i},\qquad \cos t =\frac{e^{it}+e^{-it}}2 =\frac{z+z^{-1}}2}
\end{uncoverenv}
\EMgap
\begin{EMbox}{imp}{4}{}
\begin{uncoverenv}<5->
\EMeq{\int_0^{2\pi}f(a\cos t,\,b\sin t)\,dt =\oint_Cf\!\left(a\,\frac{z+z^{-1}}{2},\;b\,\frac{z-z^{-1}}{2i}\right)\frac{dz}{iz}}
\end{uncoverenv}
\end{EMbox}
\end{frame}

\def\EMtitlesize{26}\def\EMtitleskip{31}
\begin{frame}{انتگرال‌های دستۀ دوم (ادامه)}
\begin{minipage}[t]{500bp}\vspace{0pt}
\begin{uncoverenv}<2->
\EMp{در اینجا \EMim{C} دایرۀ واحد است و به سهولت با استفاده از روش مانده‌ها انتگرال حل می‌شود.}
\end{uncoverenv}
\EMgap
\end{minipage}
\EMside{67.9bp}{\uncover<3->{\EMfigure{1.500}{m14-unit-circle}{دایرۀ واحد \EMim{z=e^{it}}، \EMim{0\le t\le2\pi}}}}
\end{frame}

\def\EMtitlesize{26}\def\EMtitleskip{31}
\begin{frame}{انتگرال‌های دستۀ دوم (ادامه)}
\begin{EMbox}{ex}{1}{}
\EMp{مطلوب است محاسبۀ انتگرال \EMim{\displaystyle\int_0^{2\pi}\frac1{2+\cos x}\,dx}.}
\EMgap
\begin{uncoverenv}<2->
\EMeq{z =e^{ix},\qquad dz =ie^{ix}dx,\qquad dx =\frac{dz}{iz},\qquad \cos x =\frac{z+z^{-1}}2 =\frac{z^2+1}{2z}}
\end{uncoverenv}
\EMgap
\begin{uncoverenv}<3->
\EMeq{2+\cos x =2+\frac{z^2+1}{2z} =\frac{z^2+4z+1}{2z}}
\end{uncoverenv}
\EMgap
\begin{uncoverenv}<4->
\EMeq{\int_0^{2\pi}\frac1{2+\cos x}\,dx =\oint_C\frac{2z}{z^2+4z+1}\,\frac{dz}{iz} =\frac2i\oint_C\frac1{z^2+4z+1}\,dz \\[4pt] =\frac2i\times2\pi i\times\sum_k\operatorname{Res}\left\{\frac1{z^2+4z+1}\right\}\Bigg|_{z=z_k}}
\end{uncoverenv}
\EMgap
\begin{uncoverenv}<5->
\EMeq{z^2+4z+1 =0\; \Longrightarrow\;z_1 =-2+\sqrt3 \approx-0.27,\qquad z_2 =-2-\sqrt3 \approx-3.73}
\end{uncoverenv}
\end{EMbox}
\end{frame}

\def\EMtitlesize{26}\def\EMtitleskip{31}
\begin{frame}{انتگرال‌های دستۀ دوم (ادامه)}
\begin{minipage}[t]{500bp}\vspace{0pt}
\begin{EMbox}{ex}{1}{}
\begin{uncoverenv}<2->
\EMp{فقط \EMim{z_1} داخل دایرۀ واحد قرار دارد.}
\vspace{3bp}
\EMeq{\int_0^{2\pi}\frac1{2+\cos x}\,dx =4\pi\times\operatorname{Res}\left\{\frac1{z^2+4z+1}\right\}\Bigg|_{z=z_1} \\[4pt] =4\pi\left[\big(z+2-\sqrt3\big)\frac1{\big(z+2-\sqrt3\big)\big(z+2+\sqrt3\big)}\right]_{z_1=-2+\sqrt3} \\[4pt] =\frac{2\pi}{\sqrt3}}
\end{uncoverenv}
\end{EMbox}
\end{minipage}
\EMside{233.6bp}{\uncover<3->{\EMfigure{1.500}{m14-poles-cos}{دایرۀ واحد و قطب‌های \EMim{z_1\approx-0.27} (داخل) و \EMim{z_2\approx-3.73} (بیرون)}}}
\end{frame}

\def\EMtitlesize{26}\def\EMtitleskip{31}
\begin{frame}{انتگرال‌های دستۀ دوم (ادامه)}
\begin{EMbox}{exe}{1}{}
\EMp{با استفاده از قضیۀ مانده‌ها مقدار انتگرال‌های معین زیر را به دست آورید:}
\vspace{3bp}
\EMeq{I_1 =\int_0^{2\pi}\frac{dt}{1-2p\cos t+p^2},\qquad I_2 =\int_{-\infty}^{\infty}\frac{\sin2x}{x^2+x+1}\,dx}
\EMgap
\begin{uncoverenv}<2->
\EMeq{I_3 =\int_{-\infty}^{\infty}\frac{x}{(x^2-2x+2)^2}\,dx,\qquad I_4 =\int_{-\infty}^{\infty}\frac{\cos^4x}{(x^2+1)(x^2+4)}\,dx}
\end{uncoverenv}
\end{EMbox}
\end{frame}
